\section{三柄球面闭包}\label{sec:s-closure}

本节记录预期的 S 论证。闭流形 HJ
定理~5.3 不用于得出 $B$ 已固定。去掉最后的 $4$-柄，称
剩余柄体为 $W_3$。其边界为 $S^3$。反向三个
$3$-柄即在 $S^3$ 上附着三个三维 $1$-柄，故
\[
\partial W_2\cong\#^3(S^1\times S^2).
\]
实际 attaching 球面系的补连通，因为在该系上球面
surgery 给出 $W_3$ 的连通边界。

\begin{lemma}[更换完整球面系下的核不变性]
\label{lem:Ssystem}
令 \(\mathcal S\) 与 \(\mathcal A\) 为
\(\#^3(S^1\times S^2)\) 中的完整球面系。若它们经 isotopy、置换与
球面 slide 相关，则自 \(\mathcal S^2_0(W_2)\) 出发的两个总 MWW 三柄映射的核
一致。
\end{lemma}

\begin{proof}
attaching 球面的 isotopy 给出同构的柄附着。
置换仅重排不相交三柄。球面 slide 恰为
\(3\)--\(3\) 柄 slide。标准柄演算给出
结果三柄配边之间的微分同胚，在其入射边界 \(W_2\) 上为恒等。Skein-lasagna
映射的自然性因此给出 \(q_{\mathcal A}=\Phi q_{\mathcal S}\)，其中
\(\Phi\) 为目标模的同构。故它们的核一致。
\end{proof}

在 Johnson 替换反向
$1$-柄图景中选取标准完整系
\(\mathcal A=(A_1,A_2,A_3)\)，已与 P0 重构立方体 $B$ 不相交。
闭流形 HJ 定理~5.3 不用于得出 $B$ 已固定。
它仅与引理~\ref{lem:Ssystem} 一起使用：任何其他完整系，
含 going-up Kirby attaching 系，在闭流形中经 isotopy、
置换与 slide 相关，故核一致。
Horvat--Jab\l{}onowski 当前预印本 \cite{HorvatJablonowski2025}，
日期 2026 年 8 月 24 日，含引理~5.5（带点子球 slide 引理）与
引理~5.7（完整球面系的相对唯一性）。本文不调用这两条引理：核不变性仅使用其定理~5.3 与
引理~\ref{lem:Ssystem}。引理~5.7 针对具球面边界流形中的相对球面系，不是本节所用设定。
在反向 $1$-柄图景中，三个 belt 球面避开 P0 立方体，
dual 圈配对为恒等，且与 C1 剩余 link 及 C2 支撑不相交。该抽象 belt 模型的 $b=0$ foam 不是
在实际二柄核心圆盘恢复之前的 MWW 曲面。
记录的球面列给出代数下界计数
\[
(b_1,b_2,b_3)=(9920,1430,311).
\]
跟踪定向对偶圆盘词而不消去非相邻相反副本给出实际几何计数
\[
(b_1,b_2,b_3)=(12578,1824,409).
\]
其乘积法向副本轮廓与相应
$\epsilon^{\otimes b_j}\Delta^{b_j-1}$ 计算显式。
相对 PL 边界映射以层级表示：它施加全部 93 个
环境对偶圆盘因子，再施加两次 Kirby 边界映射的全部六条与 1513 条带，每个所得边界副本绑定到实际 owner。
矩阵本身不再是外部输入：将对偶 handlebody 的三个正方形压缩圆盘经全部 93 个环境
Johnson 因子运输，逐因子给出 $\wedge^2A=A^{-T}$，且
mapping torus 边界为 $I-A^{-T}$。单独的曲面运输
验证器重放两次 Kirby 领圈并拒绝缺失带。

MWW 对 $3$-柄商的内在重述使用
$\mathcal S^2_0(S^2\times D^2)$ 的局部模作用。在 $N=2$ 时，令 $A_0$ 为
带一个点的本质球面，$A_1$ 为无点球面。商
关系为
\[
A_0=1,\qquad A_1=0.
\]
若 $J_j$ 为 $A_j$ 的赤道，MWW 的 K\"unneth 识别给出
两半球映射的如下精确描述：
\[
\begin{array}{ccc}
\mathcal S^2_0(W_2;J_j)&\xrightarrow{\ (\Psi_{\Delta^+},\Psi_{\Delta^-})\ }&
\mathcal S^2_0(W_2)\oplus\mathcal S^2_0(W_2)\\
\big\downarrow\wr&&\big\downarrow{\ell\oplus\ell}\\
\mathcal S^2_0(W_2)\otimes\Q\{1,X\}
&\xrightarrow{\ (\ell\circ(1\otimes\epsilon),\,\ell\circ\mu_{A_j})\ }&
\Q\oplus\Q.
\end{array}
\tag{30}
\]
此处 $\epsilon(X)=1$，$\epsilon(1)=0$，而
$\mu_{A_j}(v\otimes X)=vA_0$ 且
$\mu_{A_j}(v\otimes1)=vA_1$。因此 S 等价于两个
全源恒等式
\[
\ell(vA_0)=\ell(v),\qquad \ell(vA_1)=0
\quad\text{对每个 }v\text{ 及 }j=1,2,3.
\tag{31}
\]
这将形式映射与 MWW
定理~3.7 中的 coequalizer 对等同，但不是实际本质球面的端点求值。

\begin{proposition}[固定检测器的球面替换]\label{thm:Sgeometry}
三柄总核可用 Johnson 替换反向 $1$-柄图景的标准 belt 球面系计算，该系避开
$B$，而 $B$ 中检测器保持固定。对此系，S 约化为
六个方程 (31)。
\end{proposition}

\begin{proof}
Johnson 替换的反向 $1$-柄图景放置三个
避开 $B$ 的 belt 球面，dual 圈的柄段与
owner belt 球面相交一次，且坐标图返回避开每个 belt 立方体。
闭流形 HJ 定理~5.3 将任何其他完整
系在闭流形中经 isotopy 与此系相关；引理~\ref{lem:Ssystem}
识别核。检测器留在那些 belt
球面避开的 $B$ 中。六个方程 (31) 于是为引理~\ref{lem:Sendpoint}。
\end{proof}

\subsection{本质球面端点比较}

\begin{lemma}[标准球面的端点分解]\label{lem:Sendpoint}
对每个 \(j\)、每个缆态及每个源元 \(v\)，
\(\mu_{A_j}(v\otimes a)\) 实际端点像的常数项是 \(v\) 的旧检测器像与
穿孔标准球面秩二 TQFT 映射的张量积。因此除法三次行满足
\[
\ell(vA_0)=\ell(v),\qquad \ell(vA_1)=0.
\]
\end{lemma}

\begin{proof}
使 \(A_j\) 与二柄余核横截，令 \(b_j\) 为
核心圆盘数。MWW 定理~3.10 的证明将
\(\Delta^-\) 半球（在核心圆盘恢复之前）表示为从 \(A_j\) 删去那些圆盘及
\(\Delta^+\) 圆盘后得到的连通亏格零曲面。在一柄切割后选取 generic movie。
因 \(A_j\) 与检测器片不相交，movie 的每个临界点完全属于两曲面之一。混合事件
因此仅为可逆 Reidemeister/枢轴移动，在端点由缆 braid 表示；无混合 birth、saddle 或 death。

在 \(q=1\) 时，引理~\ref{lem:C2} 将每个混合 braid 识别为
张量范畴的对称性。birth、saddle 与 death
对该对称性的自然性将所有混合 braid 移到 movie 末端。同一端点置换与枢轴坐标图传输
检测器行与曲面映射，故它们经同时
共轭抵消。故实际常数映射为
\[
\operatorname{Id}_{\rm old}\otimes\Delta^{b_j-1}\quad(b_j>0),
\]
当 \(b_j=0\) 时为 \(\operatorname{Id}_{\rm old}\otimes\epsilon\)。
恢复的二柄核心圆盘给出 \(b_j\) 个实际 counit。故
\[
\epsilon^{\otimes b_j}\Delta^{b_j-1}(1)=0,\qquad
\epsilon^{\otimes b_j}\Delta^{b_j-1}(X)=1
\tag{32}
\]
对 \(b_j>0\)，当 \(b_j=0\) 时解释为
\(\epsilon(1)=0,\epsilon(X)=1\)。
\(1,X\) 基由 MWW
例~3.8 中 \(\Delta^+\) cap 归一化，BHPW strict 函子性防止
\(\Delta^-\) movie 中的独立符号。

完整混合 braid 映射形如 \(P(I+O(h))\)。置换
\(P\) 刚被抵消，以 \(I+O(h)\) 预复合或后复合始于
$h^3$ 阶的行仅改变到 $h^4$ 阶及以上。
故常数计算 (32) 正是除法三次行所见的候选计算。对三个非零值
$b_j=12578,1824,409$，有限 Frobenius movie 分别有
$b_j-1$ 个余积鞍点与 $b_j$ 个 counit，并再现 (32)。故
几何 $\partial W_2$ 绑定存在；下述定理在全局接口打包完整 MWW 映射之前仍停留于所述断言边界。
\end{proof}

端点交换方块是上述引理~\ref{lem:Sendpoint} 的迭代余积--counit 模型，绑定到实际层级 $W_2$ 曲面。
\[
\begin{array}{ccc}
\mathcal S^2_0(W_2;J_j)&\longrightarrow&
\mathcal S^2_0(W_2)\oplus\mathcal S^2_0(W_2)\\
\downarrow&&\downarrow\\
\mathcal S^2_0(W_2)\otimes\Q\{1,X\}&\longrightarrow&\Q\oplus\Q,
\end{array}
\]
其下映射为实际 counit 行及引理~\ref{lem:Sendpoint} 中 movie 分解的实际 \(A_j\)-作用。

\begin{theorem}[相对三柄闭包 S]\label{thm:Sdischarge}
假设~\ref{hyp:P2} 对显式 Johnson 表述成立。
\end{theorem}

\begin{proof}
三个 H1 压缩圆盘经全部 93 个 Johnson 因子与全部 1519 条 Kirby 带运输。其几何边界词长为
$12578,1824,409$；所得穿孔球面 movie 在实际 $\partial W_2$ 上实现引理~\ref{lem:Sendpoint} 的映射。
\end{proof}
